\section{결론}

본 논문에서는 동일한 LLM 기반 멀티 에이전트 환경에서 Code, LLM, Hybrid
세 가지 오케스트레이션 구조를 구현하고, 실행 시간·품질·토큰 사용량·일관성
등의 지표로 비교하였다.

실험 결과, Code 구조는 흐름이 결정론적이어서 \textbf{빠르고 안정적}이며
토큰 비용이 낮았다. LLM 구조는 흐름을 동적으로 구성하여 \textbf{유연하고
적응성}이 우수한 반면 비용과 편차가 컸다. Hybrid 구조는 두 방식의 장점을
절충하여 \textbf{가장 균형적}인 특성을 보였다.\footnote{구체적 수치는
실험 완료 후 본문 표~\ref{tab:result}에 반영한다.}

이러한 관찰을 바탕으로, 작업 특성에 따라 적합한 오케스트레이션을 선택하기
위한 설계 가이드라인을 제안하였다. 정형 반복 업무에는 Code, 창의·불확실
업무에는 LLM, 혼합형 및 대규모 운영 환경에는 Hybrid가 적합하다.

향후 연구로는 본 비교 방법론을 문서 생성(Book)에 국한하지 않고 보고서
생성(Report), 코딩 에이전트(Coding Agent), 데이터 분석(Data Analysis),
AI 어시스턴트 등 다양한 멀티 에이전트 시스템으로 확장하여 가이드라인의
일반성을 검증할 계획이다.
